\documentclass[11pt]{article}

% === Page Layout ===
\usepackage[a4paper, margin=1in]{geometry}

% === Packages ===
\usepackage{graphicx}
\usepackage{comment}
\usepackage{fancyhdr}
\usepackage{enumitem}
\usepackage{hyperref}
\usepackage{amsmath}

% === Header/Footer Setup ===
\pagestyle{fancy}
\fancyhf{}
\fancyhead[L]{Multimodal Deep Learning}
\fancyhead[R]{Answer Sheet 7}
\fancyfoot[C]{\thepage}

\begin{document}

% === Title Information ===
\begin{center}
    {\Large \textbf{Multimodal Deep Learning}} \\ [0.3em]
    {\large \textbf{Answer Sheet 7}} \\ [0.3em]
    \large Date: July 1, 2025 \\
\end{center}

\vspace{1cm}

\subsection*{Exercise 1: Selective Fine-Tuning}

\begin{enumerate}[label=(\alph*)]
    \item \href{https://arxiv.org/abs/2411.05752}{\textbf{Fishmask}} 

    \begin{enumerate}[label=(\roman*)]
        \item Mapping\\ 
        \textbf{Problem Statement:} 
        
        \begin{itemize}
            \item A layer contains \(4.0 \times 10^6\) weights.
            \item A "Fisher score approximation" assigns a scalar importance \(I_i\) to each weight \(w_i\).
            \item We are allowed to update only 0.8\% of these weights.
            \item This refers to the Fisher Information Matrix (FIM), which measures how sensitive the loss function is to parameter changes.
            \begin{itemize}
                \item High Fisher score \(I_i\): important weight.
                \item Low Fisher score \(I_i\): less important.
                \item One-shot Fisher approximation computes importance without retraining.
            \end{itemize}
        \end{itemize}
        
        \textbf{Answer:}
        
        \begin{itemize}
            \item Total weights: 4 million parameters.
            \item Budget: 0.8\% of total weights = 32,000 weights = \(k\).
            \item A binary mask \(\mathbf{m} \in \{0, 1\}^n\) has 32,000 ones and the rest zeros. (\(m_i = 1\): update; \(m_i = 0\): freeze)
            \item Mapping formula (using both \(\mathbf{I}\) and \(\mathbf{m}\)): \(\mathbf{M} = \mathbf{I}^T \mathbf{m}\)
            \item Mapping formula (constructing \(\mathbf{m}\) from top-\(k\)): \(\mathbf{m} = \mathbf{I} > \min(\operatorname{Top}_k(\mathbf{I}))\)
        \end{itemize}
    \end{enumerate}
\end{enumerate}


\end{document}